% Source PDF: source/普通高考/2011/2011浙江文.pdf, page 1, question 14.
% The source PDF contains vector line art (pdfimages -list returned no image objects).
% Node -> directed-edge list: 开始 -> k=2 -> k=k+1 -> a=4^k -> b=k^4 -> a>b?;
% a>b? --是--> 输出k -> 结束; a>b? --否--> right return arrow -> stem above k=k+1.
% All node text is locally zero-indented and centered; the return arrow points left
% into the stem and no horizontal segment carries an extra arrowhead.

\begin{tikzpicture}[
  >=Stealth,
  x=0.9cm,
  y=0.96cm,
  line width=0.45pt,
  line cap=round,
  line join=round,
  font=\normalsize,
  flowtext/.style={anchor=center, align=center,
    execute at begin node={\setlength{\parindent}{0pt}}},
  every node/.style={flowtext},
  flowbox/.style={flowtext, draw, minimum height=0.72cm,
    inner xsep=3pt, inner ysep=2pt,
    execute at begin node={\setlength{\parindent}{0pt}}},
  terminal/.style={flowbox, rounded corners=0.18cm, minimum width=1.45cm},
  process/.style={flowbox},
  decision/.style={flowbox, diamond, aspect=2.8, minimum width=3.65cm,
    minimum height=1.05cm, inner sep=0pt},
  io/.style={flowbox, trapezium, trapezium left angle=72,
    trapezium right angle=108, minimum width=2.15cm, minimum height=0.78cm},
  branchlabel/.style={flowtext, inner sep=0pt},
]
  \node[terminal] (start) at (0,12.15) {开始};
  \node[process, minimum width=1.75cm] (init) at (0,10.75) {$k=2$};
  \node[process, minimum width=2.65cm] (increment) at (0,9.20) {$k=k+1$};
  \node[process, minimum width=2.25cm] (pow4) at (0,7.65) {$a=4^k$};
  \node[process, minimum width=2.25cm] (powk) at (0,6.10) {$b=k^4$};
  \node[decision] (test) at (0,4.45) {$a>b?$};
  \node[io, minimum width=2.15cm] (output) at (0,2.55) {输出 $k$};
  \node[terminal] (finish) at (0,0.95) {结束};

  \coordinate (loop-join) at (0,10.05);
  \coordinate (loop-right) at (3.35,4.45);

  \draw[->] (start.south) -- (init.north);
  \draw[->] (init.south) -- (increment.north);
  \draw[->] (increment.south) -- (pow4.north);
  \draw[->] (pow4.south) -- (powk.north);
  \draw[->] (powk.south) -- (test.north);

  % Decision branches: only the main downward path and return edge carry arrows.
  \draw[->] (test.south) -- node[branchlabel, right=2pt] {是} (output.north);
  \draw[->] (output.south) -- (finish.north);
  \draw (test.east) -- node[branchlabel, above=2pt] {否} (loop-right)
    -- (3.35,10.05);
  \draw[->] (3.35,10.05) -- (loop-join);
  % The main stem owns the downward arrow; the return arrow ends at loop-join.
  \draw (loop-join) -- (increment.north);
\end{tikzpicture}
